네트워크 인프라의 규모 확장과 서비스 다양화로 인해 네트워크 구성 및 운영의 복잡성이 급격히 증가하고 있다. 이에 따라 운영 효율성을 향상시키고 관리 부담을 줄이기 위한 네트워크 관리 자동화 기술이 요구되고 있다. 이 과정에서 Large Language Models(LLM)을 네트워크 관리 도메인에 적용하려는 연구가 활발히 이루어지고 있다. 그렇다면 LLM은 네트워크 설정 파일을 얼마나 정확하게 해석하고 분석할 수 있는가? 기존 벤치마크는 표준 문서의 지식 회상이나 설정 생성에 치중하여 이 질문에 답하기 어렵다. 본 연구는 벤치마크, 평가 지표, 도구 통합 시스템을 하나의 프레임워크로 제안한다. NetConfigQA~2.0은 설정 파일과 Batfish 시뮬레이션을 결합한 이중 경로 파이프라인을 통해 5단계 인지 난이도의 QA 쌍을 자동 생성하며, 설정 파일과 정책 파일만 교체하면 토폴로지 규모에 관계없이 데이터셋이 확장된다. 답변 타입별 정답 구조를 반영하는 Type-Aware Accuracy(TA-Acc)를 도입하여 네트워크 데이터의 평가 문제를 다룬다. 4개 토폴로지(10\textasciitilde40 노드, 총 9,462 QA)와 6개 LLM을 평가한 결과, 사실 추출(L1\textasciitilde L3)에서는 모델 간 차이가 뚜렷하나 경로 추론과 장애 분석(L4/L5)에서는 모든 모델이 시뮬레이션 장벽에 직면하였다. 이 한계를 다루기 위해 도구 통합 Multi-Agent 시스템 NetAlly를 제시하고, Single LLM $\rightarrow$ Pure MAS $\rightarrow$ NetAlly의 3-way 비교를 통해 성능 향상의 주된 원천이 에이전트 구조가 아닌 네트워크 분석 도구의 통합에 있음을 보인다. 또한, 오류 진단, 도달 가능성 검증, 구성 합성 등의 복잡한 작업을 능동적으로 수행하도록 설계하여 실제 네트워크 운영 환경을 반영한 평가를 통해 자동화 수행 능력과 정확도를 검증하였다. 이를 통해 NetAgent가 네트워크 관리 분야에서 실질적인 적용 가능성을 가진다는 것을 입증하며, 향후 네트워크 관리 시스템 구현의 기반이 되는 기술로 활용될 수 있음을 제시한다.